\section{Related Work}
\label{sec:related}

% ../docs/RELATED_WORK.md is the repository's fuller survey. The
% characterisations in this section were verified first-hand against the full
% texts of the cited papers on 2026-08-31 (SkillOpt, the Xu-Wu SkillsBench
% study, GEPA, the coin-flip audit, both harness papers, the AGENTS.md study,
% the authoring paper, and Prompting Inversion) and on 2026-09-02 (Webson and
% Pavlick, Min et al., Ma et al., the three placebo priors in
% ../docs/RELATED_WORK.md's "What a placebo should be expected to do"); where
% this section and the docs file diverge, this section is the more recently
% checked.

\subsection{Skills as an intervention, and where the tension is}

The controlled evidence on skill files separates two questions. A study on a
30-task oracle-validated SkillsBench subset~\citep{skillsbench2026} runs two
models, six skill conditions and five trials, and splits whether a skill is
\emph{available} from how its prose is \emph{written}. Availability is a
large, robust effect: $+18.0$ to $+36.0$ percentage points over no skill,
taking the union of the two models' ranges, with every bootstrap interval
above zero. Presentation granularity, tested through content-controlled
rewrites of the same skills, is worth $+0.7$ points on one model and $-6.7$
on the other, both intervals crossing zero. The retrieval matters; the
tested rewrites of the writing do not.

SkillOpt~\citep{skillopt2026} reports $+23.5$ points from optimising exactly
that prose. The figure is one harness, GPT-5.5 in direct chat; the same
sentence gives $+24.8$ inside Codex and $+19.1$ inside Claude Code, and
quoting the first without the other two is the error
\citet{harnessdisclosure2026} exists to name. The two results are in tension
at exactly one point: the granularity null
says rewriting a skill's prose while holding its content fixed moves little,
and SkillOpt's gains come from rewrites that are free to change content,
selection and emphasis at once. An optimiser that earns $+23.5$ points is
therefore either adding content or exploiting something no
presentation-controlled rewrite can reach, and neither paper's design can say
which. Our design holds the model and corpus fixed, adds the control both
lack, and tests what the optimisers actually hand back.

\paragraph{SkillOpt's method, and the narrow criticism.} SkillOpt makes
bounded add, delete and replace edits to a single skill document, gates each
edit on a selection split disjoint from the rollouts the edit was generated
from, and accepts on strictly-better; committed edits per run number 11 to
44. Its harness choice is the right one: it evaluates inside the agent tools
people actually run. The criticism here is narrow. Every acceptance is a
single unreplicated evaluation, the same selection items are reused for
every decision in the run, and no confidence interval, significance test or
multiplicity correction appears anywhere in the paper: every reported result
is a point estimate from a single run on fixed splits. None of its baselines
is a document of matched length whose content says nothing. We reproduce its
engine faithfully and test its output.

\paragraph{GEPA.} GEPA~\citep{gepa2026} evolves prompts by sampling
rollouts, reflecting on successes and failures in natural language, and
proposing wholesale revisions of a module's instruction, with candidate
diversity maintained by a per-instance Pareto frontier. It reports
outperforming a reinforcement-learning baseline by 6\% on average and by up
to 20\%, at up to $35\times$ fewer rollouts. Its acceptance has the same
shape as SkillOpt's: a candidate enters the pool when its score on a
reflection minibatch, three items by default, beats its parent's, and no
significance notion appears in the algorithm or the evaluation. We use it as
distributed, through the library's documented adapter interface, with an
OpenAI-compatible reflection model; the skill here is a single document, so
a module revision is a whole-prompt revision.

\subsection{Why a placebo arm is the load-bearing control}

Presence-versus-absence designs dominate this literature. The AGENTS.md
impact study~\citep{agentsmd2026} compares repositories with and without the
file and reports efficiency gains; its two conditions are the file present
and the file removed. The skill-authoring
literature~\citep{authoringskills2026} prescribes how to write skills and
reports no comparisons at all. The SkillsBench benchmark
itself~\citep{skillsbench2026li} measures curated skills against a no-skill
baseline, and its authors flag that the observed gains could partly reflect
added context length, calling for length-matched baselines of random or
irrelevant text as future work. In a design without that baseline, ``the
skill helped'' and ``a document-shaped block of tokens in the system prompt
helped'' are the same observation.

What a document-shaped placebo should be expected to do is not new to skill
evaluation. Webson and Pavlick find that NLI models learn as fast from
templates that are intentionally irrelevant or misleading as from
instructive ones in the few-shot setting, while a null template with no text
at all does worse than an irrelevant
one~\citep{webson2022understand}. Min et al. find that scrambling which
label goes with which example in an in-context demonstration, while keeping
its label set, input distribution and format, costs little, and that
removing the format instead costs as much as dropping the demonstration
entirely; the label set and the input distribution are content by the
paper's own accounting, so we read this for the correctness of the pairing
and not for content in general~\citep{min2022demonstrations}. Ma et al. find
that a reflective prompt optimiser's feedback is much the same whether it
reflects on real errors or on errors fabricated by flipping predictions, and
that a directionless resampler matches its gains; the paper does not say
where those gains come from, and we do not read it that
way~\citep{ma2024promptoptimizers}. Together the three say that a block of
structured text with no method in it is expected to move a model, by an
amount that tracks its shape more than its content, in settings adjacent to
ours: a skill that beats an empty prompt has shown less than it appears to
until it also beats a placebo built to the same shape.

Our placebo is generated to match the shipped skill on word count and on
structural shape, and it ships in the repository beside the skill so the
comparison can be re-run. \Cref{sec:results} reports it as a first-class
arm, every arm in the registered family is tested against it, and
\Cref{sec:length-confound} reports what it does and does not control for the
evolved winners.

\subsection{Harness variance, and what a verdict is worth}

\citet{harnessdisclosure2026} run three frontier models across three harness
configurations on a stratified subset of SWE-bench Verified and report an
aggregate harness-to-model variance ratio of $7.8\times$, with the ranking
of models reversing in six of nine comparisons. One $3\times3$ design on one
task distribution: we read the direction as transferring and the ratio as
not. \citet{harnessbench2026} reach the qualitative half of the same
conclusion over 5{,}194 trajectories on 106 sandboxed tasks, where
completion, process quality and failure behaviour all vary substantially
across model-harness pairings. A result measured inside one scaffold is a
claim about that scaffold as much as about the model. We give a full
disclosure in \Cref{sec:appendix-harness}, against the seven-layer
requirement \citet{harnessdisclosure2026} sets, and we do not assume our
result transfers to a different agent harness.

Prompting Inversion~\citep{promptinginversion2025} sharpens this for our
case: prompting strategies that helped a mid-tier model became detrimental
on a frontier one. Our engines optimised against a 1.7B model. That is a
reason to expect the transfer question to have its own answer, and a reason
we decline to extrapolate from ours.

\subsection{Statistics for evaluations at this scale}

\citet{miller2024errorbars} popularised CLT-based standard errors for
evaluations. \citet{noclt2025} shows the normal approximation is unreliable
below a few hundred effectively independent datapoints, which is where a
clustered design at our item count sits. We use McNemar's exact
test~\citep{mcnemar1947note} for paired binary outcomes and a cluster
bootstrap over templates for intervals, and we report the discordant split
alongside every $p$-value.

Audited head-to-head, prompt optimisation itself is close to a coin flip:
across six optimisers, none of them GEPA, on four single-agent tasks under
equal compute, 49\% of 72 optimisation runs scored below
zero-shot~\citep{promptcoinflip2026}. The multiplicity problem in this
neighbourhood is larger than the reported comparisons suggest, because the
ratchet inside each engine performs comparisons that never reach a table. We
cannot correct for those. We correct for ours, across a family of three per
item set fixed before the first call, with Holm's step-down procedure.

Benjamini--Hochberg~\citep{benjamini1995fdr} is implemented in the same
module and controls the false discovery rate rather than the family-wise
error rate; with no rejection at either criterion, the choice does not
change the result.
